\section{Signal Strength} \label{sec:signal_strength}
The RF performance of the custom PCB was evaluated by measuring the RSSI at varying temperature and distance between the end device and the Zigbee coordinator. RSSI is reported in \si{\dBm}, a logarithmic measure of power referenced to \SI{1}{\milli\watt}, and represents the received power at the antenna port of the receiver.
The EFR32MG22E transceiver reports RSSI with a resolution of \SI{0.25}{dBm} and an accuracy of \SI{\pm 6}{\decibel} in the linear operating region, as defined by the datasheet~\cite{EFR32MG22E}.
The relevant lower bound for link evaluation is the receiver sensitivity, which is specified at \SI{-102.3}{dBm}. Any measured RSSI above this threshold confirms that the received signal is within the demodulation capability of the hardware.
As a practical operational reference, the Institute of Electrical and Electronics Engineers (IEEE) 802.15.4-2024 standard specifies a minimum receiver sensitivity requirement of \SI{-85}{dBm} for compliant devices, which represents a conservative link margin target under which reliable communication is expected across a broad range of deployments~\cite{IEEE8021542024}.

In a full-system test in the climate chamber, the RSSI was measured for a transmit power of \SI{3}{\dBm}, with the receiver placed outside the chamber behind two \SI{5}{\centi\meter} rubber insulation sheets. The result is presented in \Cref{fig:FS_RSSI} and shows a mean of \SI{-47}{\dBm}, with a minimum of \SI{-54}{\dBm} and a maximum of \SI{-43}{\dBm}. The standard deviation across the full \SIrange{-20}{50}{\celsius} sweep is only \SI{1.8}{\decibel}, which is well within the \SI{\pm6}{\decibel} RSSI accuracy of the transceiver. The link strength therefore shows no clear dependence on temperature, the observed variation being comparable to the measurement uncertainty.

\begin{figure}
\centering
\includestandalone[mode=buildmissing]{Standalone/Plots/FullSystem_neg_deg_RF}
    \caption{Received signal strength (RSSI) as a function of temperature over the \SIrange{-20}{50}{\celsius} range at \SI{30}{\percent} relative humidity.}
    \label{fig:FS_RSSI}
\end{figure}

A distance test was performed to measure the RSSI of both the encapsulated and the non-encapsulated module. The test could not be conducted in free air, as a connection to the WiFi network was required, and was therefore carried out along a long hallway in an office building. Each prototype was measured from \SI{0}{\meter} to \SI{30}{\meter} in \SI{5}{\meter} intervals, at three transmit power levels of \SI{0}{\dBm}, \SI{3}{\dBm} and \SI{6}{\dBm}. The results are shown in \Cref{fig:DistanceTest}.

\begin{figure}
\centering
\includestandalone[mode=buildmissing]{Standalone/Plots/signalStrength}
    \caption{Signal strength (RSSI) as a function of module distance from the coordinator.}
    \label{fig:DistanceTest}
\end{figure}
 
Across all distances and power levels, the encapsulated module reads consistently weaker than the bare module, by an average of approximately \SI{3}{\decibel}, which is attributed to attenuation of the signal by the silicone encapsulant. Taking the reference sensitivity of \SI{-85}{\dBm} as the minimum acceptable received level, this loss reduces the usable range by one \SI{5}{\meter} interval at each power setting, as summarized in \Cref{tab:DistanceRange}. The bare module maintains an acceptable link out to \SI{15}{\meter}, \SI{20}{\meter} and \SI{25}{\meter} at \SI{0}{\dBm}, \SI{3}{\dBm} and \SI{6}{\dBm} respectively, while the encapsulated module reaches \SI{10}{\meter}, \SI{15}{\meter} and \SI{20}{\meter} at the same settings. Each \SI{3}{\dBm} increase in transmit power recovers one interval of range, so the attenuation introduced by the encapsulant can be fully compensated by operating one power step higher. At the maximum tested power of \SI{6}{\dBm}, the encapsulated module sustains a reliable link out to \SI{20}{\meter}.

\begin{table}
    \centering
    \caption{Maximum range at which the received signal remains at or above the \SI{-85}{\dBm} acceptability threshold, for the bare and encapsulated modules at each transmit power.}
    \label{tab:DistanceRange}
    \begin{tabular}{ccc}
        \toprule
        Transmit power [\si{\dBm}] & Bare module [\si{\meter}] & Encapsulated module [\si{\meter}] \\
        \midrule
        0 & 15 & 10 \\
        3 & 20 & 15 \\
        6 & 25 & 20 \\
        \bottomrule
    \end{tabular}
\end{table}
 

% Normal PCB (no silicone)							
%  	0m	5m	10m	15m	20m	25m	30m
% 0 dBm	-33	-70	-81	-85	-92	-94	-101
% 3 dBm	-30	-66	-77	-80	-84	-89	-95
% 6dBm	-25	-63	-73	-76	-80	-84	-90

% Encapsulated PCB (with silicone)							
% 	0m	5m	10m	15m	20m	25m	30m
% 0 dBm	-36	-77	-82	-86	-93	-96	-104
% 3 dBm	-32	-70	-80	-82	-89	-94	-99
% 6dBm	-28	-67	-75	-79	-85	-89	-93


\section{Integrated Temperature Sensor} \label{sec:temp_sensor}

The strain sensors have demonstrated temperature sensitivity when the Wheatstone bridge is not fully balanced, as discussed in \Cref{Temperature Behaviour}. A potential mitigation strategy is to mathematically compensate the measured sensor output as a function of temperature, for which a reliable on-board temperature reference is required. The on-die temperature sensor of the EFR32MG22E is therefore evaluated to determine its suitability for this purpose. The sensor is embedded within the silicon die of the MCU and reports absolute temperature with a nominal resolution of \SI{0.25}{\celsius} and a
specified single-measurement noise floor of \SI{0.6}{\celsius}\,RMS~\cite{EFR32MG22E}.
The datasheet further specifies a typical mean temperature offset of \SI{3.14}{\celsius}
relative to ambient, attributable to self-heating of the die under normal operating
conditions.

To characterize the temperature sensor's performance, the PCB was subjected to a controlled thermal profile in a climate chamber spanning \SIrange{-20}{50}{\celsius} at a constant relative humidity of \SI{30}{\percent}.
\Cref{fig:FS_temp} shows the chamber reference temperature alongside the on-die sensor reading over the full test duration.

\begin{figure}
    \centering
    \includestandalone[mode=buildmissing]{Standalone/Plots/FullSystem_neg_deg_temp}
    \caption{On-die temperature sensor measurement under varying environmental conditions
             (\SIrange{-20}{50}{\celsius}, \SI{30}{\percent} relative humidity).}
    \label{fig:FS_temp}
\end{figure}

\begin{figure}
    \centering
    \includestandalone[mode=buildmissing]{Standalone/Plots/MCUtemp_vs_AmbientTemp}
    \caption{Temperature offset, defined as the difference between the on-die and the ambient reference temperature, over the full climate-chamber test.}
    \label{fig:temp_diff}
\end{figure}

The mean signed deviation between the on-die measurement and the ambient reference was found to be \SI{+3.42}{\celsius}, as illustrated in \Cref{fig:temp_diff}, consistent with the datasheet-specified offset.
Of greater significance for the intended compensation scheme is the standard deviation of the offset, which was found to be $\sigma = \SI{1.89}{\celsius}$ across the full test profile.
A purely static offset would manifest as a near-zero standard deviation, so the elevated value observed over the complete profile is attributed to a thermal lag effect, whereby the thermal mass of the IC package introduces a delay in the die temperature response relative to changes in the ambient temperature, causing the offset to vary systematically during temperature ramps.

\begin{figure}
    \centering
    \includestandalone[mode=buildmissing]{Standalone/Plots/MCUtemp_vs_AmbientTemp_correlation}
    \caption{Scatter plot of the on-die measured temperature against the ambient reference
             temperature over the full climate-chamber test.}
    \label{fig:temp_scatter}
\end{figure}

The relationship between the on-die measured temperature and the chamber reference temperature is shown as a scatter plot in \Cref{fig:temp_scatter}. A linear regression yields
\begin{equation}
    T_{\mathrm{MCU}} = 0.978\,T_{\mathrm{amb}} + \SI{3.92}{\celsius}
    \label{eq:temp_fit}
\end{equation}
with a coefficient of determination $R^2 = 0.994$, indicating an excellent linear fit throughout the tested temperature range. The slope of $0.978$ deviates from unity by only \SI{2.2}{\percent}, and the intercept of \SI{3.92}{\celsius} represents the predicted on-die reading when the ambient temperature is \SI{0}{\celsius} and captures the combined effect of the systematic self-heating offset and the slope deviation from unity. %The high linearity described by $R^2$ confirms a strong linear relationship between the two measurements throughout the tested temperature range.

The largest deviations are observed during temperature ramps, as the die temperature lags
behind the chamber air temperature due to the thermal mass of the package.
To isolate the sensor's steady-state noise performance, the standard deviation of the
on-die reading was evaluated separately during periods in which the chamber temperature
was held constant.
The results are summarised in \Cref{tab:std_table}.
Under thermally settled conditions, the standard deviation ranged from
\SIrange{0.66}{0.80}{\celsius}, which is in close agreement with the specified
single-measurement noise of \SI{0.6}{\celsius}\,RMS~\cite{EFR32MG22E}.
The slight excess above the datasheet figure is expected, as the specification represents
a best-case value obtained under tightly controlled conditions.
Noise can be reduced further through the hardware averaging function provided by the EMU
peripheral, which supports 16- and 64-sample burst averaging and reduces the RMS noise to
\SI{0.17}{\celsius} and \SI{0.12}{\celsius}, respectively~\cite{EFR32MG22E}.
On this basis, the on-die temperature sensor is considered adequate for the proposed
compensation scheme, provided that corrections are applied under thermally settled
conditions.

\begin{table}
\centering
\caption{Standard deviation of the on-die temperature reading during periods of constant chamber temperature.}
\label{tab:std_table}
\begin{tabular}{S[table-format=+2.0] S[table-format=2.1] S[table-format=1.2]}
\toprule
{Chamber Temp. (\si{\celsius})} & {Time Interval (h)} & {Std. Dev. $\sigma$ (\si{\celsius})} \\
\midrule
+25 & {0.0 -- 2.0}   & 0.66 \\
+25 & {17.5 -- 19.5} & 0.78 \\
+50 & {3.3 -- 5.3}   & 0.80 \\
+50 & {14.3 -- 16.3} & 0.80 \\
-20 & {8.9 -- 10.9}  & 0.68 \\
\bottomrule
\end{tabular}
\end{table}

\subsection{Current Consumption Analysis}
\cref{fig:cc_full} shows the measured current consumption of the full system 
across a temperature range of \SI{-20}{\celsius} to \SI{+50}{\celsius} at constant 
humidity. The two operating modes exhibit opposing temperature dependencies, each 
governed by distinct physical mechanisms.

\begin{figure}[H]
\centering
\begin{subfigure}{0.48\textwidth}
    \centering
    \includestandalone[mode=buildmissing, width=\textwidth]{Standalone/Plots/FullSystem_neg_deg_cc_meas}
    \caption{EM0 (run mode)}
    \label{fig:measurement}
\end{subfigure}
\hfill
\begin{subfigure}{0.48\textwidth}
    \centering
    \includestandalone[mode=buildmissing, width=\textwidth]{Standalone/Plots/FullSystem_neg_deg_cc_sleep}
    \caption{EM4 (sleep mode)}
    \label{fig:sleep}
\end{subfigure}
\caption{Full system current consumption of EM0 (run mode) and EM4 (sleep mode).}
\label{fig:cc_full}
\end{figure}

In EM0 (run mode), the average current consumption increases at lower temperatures,  from \SI{5.612}{\milli\ampere} at \SI{+50}{\celsius} to \SI{6.681}{\milli\ampere} at \SI{-20}{\celsius}, with a total average at \SI{5.886}{\milli\ampere} for the whole time period.This trend is opposite to the intrinsic MCU behavior characterized in the EFR32MG22E datasheet \cite{EFR32MG22E}, which shows EM0 supply current increasing with temperature when powered from a stable regulated supply. The reason for this reversal can be found in the temperature dependence of the EVE ER14250 Li-SOCl\textsubscript{2} battery. The output of the battery was measured as shown in \cref{fig:FS_bat}. Here the battery terminal voltage drops significantly at low temperatures with active-pulse voltages approaching \SI{2.6}{\volt} at \SI{-20}{\celsius}. 

\begin{figure}[H]
\centering
\includestandalone[mode=buildmissing, width=0.75\textwidth]{Standalone/Plots/FullSystem_neg_deg_bat}
    \caption{Battery voltage in varying enviromental conditions.}
    \label{fig:FS_bat}
\end{figure}

Since the MCU digital core, RF subsystem, and I/O supply rails (DVDD, RFVDD, PAVDD, AVDD, and IOVDD) are all powered by the integrated DC--DC buck converter operating at \(V_{\mathrm{DCDC}}=\SI{1.8}{\volt}\), the input current can be estimated as a function of the supply voltage under the assumption of a constant load. To investigate whether the observed decrease in input current with increasing supply voltage can be explained by the DC--DC converter, the run-mode load is modeled as a constant power demand \(P_{\mathrm{out}}\) supplied by the converter. The converter efficiency is assumed to be \(\eta \approx 0.91\), based on the efficiency characteristic provided in the datasheet \cite{EFR32MG22E}. The efficiency of a buck converter is defined as the ratio of output power to input power \cite[p.~93]{erickson2020fundamentals}:

\begin{equation}
    \eta = \frac{P_{\mathrm{out}}}{P_{\mathrm{in}}}.
\end{equation}

Using that the converter input power is given by \(P_{\mathrm{in}} = V_{\mathrm{in}} I_{\mathrm{in}}\) and the \SI{50}{\celsius} operating point (\(V_{\mathrm{in},50^\circ\mathrm{C}} \approx \SI{3.4}{\volt}\), \(I_{\mathrm{in},50^\circ\mathrm{C}} = \SI{5.612}{\milli\ampere}\)) as a reference, the delivered output power is

\begin{equation}
    P_{\mathrm{out}}
    = \eta V_{\mathrm{in},50^\circ\mathrm{C}} I_{\mathrm{in},50^\circ\mathrm{C}}
    = 0.91 \cdot \SI{3.4}{\volt} \cdot \SI{5.612}{\milli\ampere}
    = \SI{17.36}{\milli\watt}.
\end{equation}

Assuming that \(P_{\mathrm{out}}\) remains approximately constant with temperature, the reduced battery terminal voltage at \SI{-20}{\celsius} (\(V_{\mathrm{in},-20^\circ\mathrm{C}} \approx \SI{2.7}{\volt}\)) yields a predicted input current of

\begin{equation}
    I_{\mathrm{in},-20^\circ\mathrm{C}}
    = \frac{P_{\mathrm{out}}}{\eta V_{\mathrm{in},-20^\circ\mathrm{C}}}
    = \frac{\SI{17.36}{\milli\watt}}
           {0.91 \cdot \SI{2.7}{\volt}}
    \approx \SI{7.07}{\milli\ampere}.
\end{equation}

The predicted current of \SI{7.07}{\milli\ampere} is in good agreement with the measured value of \SI{6.681}{\milli\ampere}, corresponding to a deviation of approximately \(5.8\%\). This small discrepancy is expected, as the estimation assumes constant converter efficiency and a temperature-independent load power, while neglecting secondary effects such as temperature-dependent converter losses, variations in RF and digital subsystem power consumption, and battery impedance effects. Nevertheless, the result supports the hypothesis that the increase in measured input current at low temperatures is primarily caused by the reduction in battery terminal voltage, requiring the DC--DC converter to draw more current to maintain approximately constant output power.

In EM4 (sleep mode), the average current consumption rises with temperature, from \SI{171.9}{\nano\ampere} at \SI{-20}{\celsius} to \SI{767.7}{\nano\ampere} at \SI{+50}{\celsius}, with a total average of \SI{451}{\nano\ampere} for the whole time period. This trend matches the EFR32MG22E datasheet and confirms that sleep current dominates the long-term energy budget at low sample rates and elevated temperatures, as reflected in \autoref{tab:charge_10y}. It also proves that the microcontroller is operating at the lower power consumption possible, since that datasheet specified the absolute max at \SI{170}{\nano\ampere}. 

% The measured average current consumption in both operating modes at \SI{-20}{\celsius} and \SI{+50}{\celsius} is presented in \autoref{tab:cc_averages}. The results are consistent with the datasheet specifications, with the lowest measured EM4 current consumption being \SI{0.17}{\micro\ampere} and the EM0 current consumption ranging from \SI{5.612}{\micro\ampere} to \SI{6.674}{\micro\ampere} over the tested temperature range.

% \begin{table}[H]
% \centering
% \caption{Measured current consumption and charge per wake-up}
% \label{tab:cc_averages}
% \begin{tabular}{@{}lcr@{}}
% \toprule
% Mode &
% Temp. &
% $I_{\mathrm{avg}}$  \\
% &
% (\si{\celsius}) &
% (\si{\milli\ampere}) \\
% \midrule
% EM0 & -20 & \SI{6.674}{\milli\ampere} \\
% EM0 & +50 & \SI{5.612}{\milli\ampere} \\
% EM4 & -20 & \SI{171.9}{\nano\ampere} \\
% EM4 & +50 & \SI{767.7}{\nano\ampere} \\
% \bottomrule
% \end{tabular}
% \end{table}

A current consumption analysis is carried out to determine the nominal battery capacity required for a given number of reports over the device lifetime. Taking a \num{10}-year lifetime, the total operating time is
\begin{equation}
    H_{10\text{y}} = 24~\text{h/day} \times 365~\text{days/year}
        \times 10~\text{years} = \SI{87600}{\hour},
\end{equation}

The charge associated with a given current consumption is obtained from
\begin{equation}
    Q = I \cdot t,
\end{equation}
where $I$ is the current and $t$ the duration over which it is drawn.

In \cref{tab:charge_10y}, the average charge per cycle is used to determine the total nominal charge required to support different numbers of readings per day over a 10-year period. For the selected EVE ER14250 \SI{1200}{\milli\ampere\hour} battery, assuming a 10-year lifetime and an estimated capacity degradation of 1\% per year (resulting in a total loss of approximately 12\%, or \SI{144}{\milli\ampere\hour}), the maximum number of cycles per day is calculated to be approximately 55 at the highest current consumption at \SI{50}{\celsius}.

\begin{table}[H]
\centering
\caption{Estimated charge consumption over 10 years}
\label{tab:charge_10y}
\begin{tabular}{@{}crrrr@{}}
\toprule
Readings/day &
Temp. &
$Q_{\mathrm{EM0}}$ &
$Q_{\mathrm{EM4}}$ &
$Q_{\mathrm{total}}$ \\
Over 10 years  &
(\si{\celsius}) &
(\si{\milli\ampere\hour}) &
(\si{\milli\ampere\hour}) &
(\si{\milli\ampere\hour}) \\
\midrule
1 & -20 & 15.24 & 15.05 & 30.29 \\
1 & +50 & 18.12 & 67.23 & 85.35 \\
\hline
2 & -20 & 30.48 & 15.05 & 45.52 \\
2 & +50 & 36.24 & 67.21 & 103.45 \\
\hline
5 & -20 & 76.19 & 15.04 & 91.22 \\
5 & +50 & 90.61 & 67.15 & 157.75 \\
\hline
10 & -20 & 152.38 & 15.01 & 167.39 \\
10 & +50 & 181.21 & 67.04 & 248.25 \\
\hline
55 & -20 & 838.07 & 14.80 & 852.87 \\
55 & +50 & 996.67 & 66.10 & 1062.77 \\
\bottomrule
\end{tabular}
\end{table}

% In summary, based on the measured current consumption of the complete system implementation, it was demonstrated that the battery input voltage has the dominant influence on EM0 (run mode) current consumption. The experimental results follow the same overall trend as reported in the \cite{EFR32MG22E} datasheet, indicating good agreement between the measured behavior and the expected device characteristics.

